\chapter{Análisis del problema}\label{chp:analysis}

En este capítulo se tratarán los requisitos \ref{sec:requirements} se ajustan al proyecto necesarios para el correcto funcionamiento del mismo.

\section{Requisitos}\label{sec:requirements}
La tabla de requisitos \ref{table:requirements} describe 4 tipos: \textbf{información} \ref{sub:information-requirement}, \textbf{funcional} \ref{sub:functional-requirement}, \textbf{no funcional} \ref{sub:non-functional-requirement} y \textbf{hardware} \ref{sub:hardware-requirement}. Cada elemento consta de un identificador, un número y una descripción.
\begin{table}[htp]
\centering

    % Esta primera línea define las columnas de la tabla. Los posibles tipos de columna son:
    % c: texto centrado
    % l: texto alineado a la izquierda
    % r: texto centrado a la derecha
    % p: columna de ancho fijo
    % Las columnas tienen ancho dinámico según la anchura máxima de los elementos que contengan.

    % Las columnas l/r/c no parten el texto en filas diferentes si éste es demasiado largo. Para ello, puede utilizar el tipo de columna de ancho fijo "p".
    
    % Las barras verticales | se usan para definir los bordes verticales de la tabla. Pruebe a eliminar algunas y observe qué ocurre.
    \begin{tabular}{ | p{3cm} | p{3cm} | p{3cm}| p{3cm} | }
        
        % A continuación van las filas de la tabla. En cada fila, las columnas se separan con el carácter &
        % Para terminar una fila se usa \\
        % Para incluir un borde horizontal entre filas se usa \hline

        % Cabecera con textos en negrita:
        \hline
        \textbf{Información} & \textbf{Funcionales} & \textbf{No funcionales} & \textbf{Hardware}\\
        \hline
        
        % Cuerpo de la tabla:
        I.01 Código GLSL & F.01 Biblioteca de shaders & NF.01 Robustez y detección de errores & H.01 Poseer una GPU configurada correctamente \\
        \hline
        I.02 Resultados, análisis y gráficas & F.02 Biblioteca de errores & NF.02 Modular y organizado & H.02 Sin mínima exigencia de CPU, RAM, disco, etc \\
        \hline
         & F.03 Biblioteca de polígonos & NF.03 Código comentado y en inglés & \\
        \hline
         & F.04 Biblioteca miscelánea de C & & \\
        \hline
         & F.05 Biblioteca de animación & & \\
        \hline
         & F.06 Ejemplo de uso y pruebas & & \\
        \hline
        
    \end{tabular} 
    
    \caption{Tabla de requisitos}
    \label{table:requirements} 
\end{table}


\subsection{Requisitos de información}\label{sub:information-requirement}
Los requisitos de información describen los documentos y código que debe ser almacenado, leído y procesado.

\begin{enumerate}

    \item \textit{Shaders}, código \textit{GLSL} externo que necesita ser compilado dentro del programa para renderizar la geometría.

    \item Resultados de los programas, almacenados en diferentes formatos con el objeto de llevar a cabo análisis de resultados y transformarlos en gráficas.
    
\end{enumerate}

\subsection{Requisitos funcionales}\label{sub:functional-requirement}
Los requisitos funcionales describen cómo se debe comportar el sistema al ser ejecutado y qué propiedades tiene.

\begin{enumerate}

    \item Debe existir una biblioteca relacionada con \textit{shaders} que facilite su uso para el usuario.

    \item Debe existir una biblioteca de errores que provea al usuario de una manera sencilla de comprobar la validez de ciertas acciones y actuar en consecuencia.

    \item Debe existir una biblioteca relacionada con \textbf{polígonos}, que construya geometría para el usuario.

    \item Debe existir una biblioteca miscelánea que cubra propósitos generales de programación, como leer archivos y crear estructuras de memoria.

    \item Debe existir una biblioteca de animación, que soporte crear una secuencia de transformaciones sobre los vértices en pantalla y que simulen movimiento, pudiendo configurar los fotogramas por segundo, acciones y demás.

    \item Deben existir programas que utilicen todas las bibliotecas creadas, para servir de ejemplo y sujeto de pruebas para los resultados a analizar.

\end{enumerate}

\subsection{Requisitos no funcionales}\label{sub:non-functional-requirement}
Los requisitos no funcionales describen características más abstractas a las funcionables que describen propiedades ajenas al comportamiento técnico del sistema, como pueden ser la robustez, la modularidad o la fluidez.

\begin{enumerate}

    \item El sistema debe ser robusto y detectar errores tras operaciones cruciales en la secuencia de los programas.
    
    \item El proyecto debe estar organizado de forma ordenada y modular.

    \item El código del proyecto debe estar comentado correctamente, además de estar escrito en el idioma inglés.
    
\end{enumerate}

\subsection{Requisitos hardware} \label{sub:hardware-requirement}

\begin{enumerate}
    \item Poseer una \textit{GPU}, preferiblemente \textit{AMD} o \textit{NVIDIA}, configurada correctamente con los controladores apropiados.

    \item No es necesario establecer un requisito mínimo de otros componentes, ya sea \textit{RAM}, \textit{CPU} o disco, pues \textit{OpenGL} es flexible.
\end{enumerate}

